\chapter{Methodology}

\section{System requirements}

\subsection{Functional requirements}
\begin{itemize}
    \item Generate a sufficiently large, high-quality, correctly answered set of teacher CoT samples for a GSM8K training subset.
    \item Persist per-sample metadata (latency, token counts, correctness, cost) for analysis.
    \item Fine-tune a single small student model to emit coherent step-by-step reasoning and a correctly formatted final answer line.
    \item Fine-tune a single small student model to emit only the correctly formatted final answer line.
    \item Evaluate all models (base, students, teacher) on the GSM8K test split under strict numeric exact-match parsing of the final answer line.
\end{itemize}

\subsection{Non-functional requirements}
\begin{itemize}
    \item \textbf{Optimized asynchronous generation}: Highly asynchronous dataset generator with concurrency tuned to balance throughput and timeout/error rate.
    \item \textbf{Observability}: Persistent logging of token usage, latency distributions, and derived cost metrics during dataset generation.
    \item \textbf{Data quality}: Diverse and coherent chain-of-thought reasoning. Removal of outliers.
    \item \textbf{Resource efficiency}: All training and benchmarking runs fit in available VRAM without OOM.
    \item \textbf{Format compliance}: Final answer line must follow ``Final Answer: <integer>'' exactly. Deviation counts as failure.
    \item \textbf{Maintainability}: Separation of data generation, cleaning, training, and evaluation stages with clear artifact boundaries.
    \item \textbf{Reproducibility}: Deterministic training pipeline, versioned artifacts, prompts, configuration, environment and scripts.
    \item \textbf{Budget adherence}: Monitor spend so teacher generation and GPU usage stay within the stated cap.
\end{itemize}

\subsection{Constraints}
\begin{itemize}
    \item Single-GPU environment - no multi-node or distributed training.
    \item Fixed public benchmark - GSM8K.
    \item Teacher accessed only via external API.
    \item Limited sequence context - excessive rationales must be pruned or truncated without losing the final answer line.
    \item Budget cap at 200\,PLN.
\end{itemize}

\subsection{Success criteria}
\begin{itemize}
    \item Accuracy improvement of the distilled student over the base CoT-prompted model by at least 5 percentage points on GSM8K test.
    \item Cleaned dataset size sufficient to enable measurable gains (tens of thousands of rationale+answer pairs) without exceeding resource limits.
    \item All runs complete within available GPU memory and budget envelope.
\end{itemize}

\subsection{Risks and mitigation}
\begin{itemize}
    \item \textbf{Overfitting to limited subset}: Mitigate via data augmentation and training regularization.
    \item \textbf{Format consistency}: Enforce prompt templates and validation during evaluation.
    \item \textbf{Data noise leakage}: Strict filtering to remove incorrect teacher answers before student training.
    \item \textbf{Cost escalation}: Concurrency and request budgeting monitored. Setting account balance limits for GPU usage.
\end{itemize}

\section{Architecture}
The system is an offline teacher-to-student pipeline that produces a cleaned CoT dataset and trains a 3B-class student with PEFT, then evaluates under strict numeric exact match.

\subsection{Components}
\begin{itemize}
    \item \textbf{Teacher generator}: Async OpenRouter client producing structured CoTs.
    \item \textbf{Prompts/config}: Versioned templates under \texttt{prompts/} with a meta JSON capturing decoding + generation config.
    \item \textbf{Dataset artifacts}: Raw, cleaned, and processed JSONL stages (introduced once in Data Flow).
    \item \textbf{Cleaning \& preprocessing}: Notebooks applying correctness filtering, rationale dedup, and length outlier removal.
    \item \textbf{Training}: Notebook performing QLoRA fine-tuning of two modes (SCoTD and label-only).
    \item \textbf{Evaluation}: Notebook for greedy exact-match benchmarking.
    \item \textbf{Runtime}: Single Python environment.
\end{itemize}

\subsection{Data flow}
\begin{enumerate}
    \item \textbf{Ingest}: Load GSM8K train/test splits; select first $\sim$1{,}000 train questions.
    \item \textbf{Generate}: Oversample $\sim$30 CoTs per question asynchronously and write raw JSONL.
    \item \textbf{Clean}: Apply correctness filter, rationale dedup, and 99th percentile length trimming and produce cleaned JSONL.
    \item \textbf{Process}: Project cleaned dataset to minimal supervised learning schema and produce processed JSONL.
    \item \textbf{Train}: Two modes (SCoTD, label-only).
    \item \textbf{Evaluate}: Strict numeric exact-match benchmarking.
\end{enumerate}

\section{Teacher CoT dataset generation}
This stage produces a cleaned high-quality rationale+answer corpus from a strong reasoning model.

\subsection{Subset selection}
For cost control only the first $\sim$1{,}000 GSM8K train questions are used. Each question targets 30 samples (oversampling for rationale diversity). Generation is resumable: existing \texttt{(question\_id, sample\_id)} pairs found in the output JSONL are skipped on subsequent runs.

\subsection{Async generation}
An async Python client (HTTPX) calls the OpenRouter \texttt{/chat/completions} endpoint for model \texttt{deepseek/deepseek-r1-distill-qwen-32b}. Concurrency is capped at 25 via an \texttt{asyncio.Semaphore}. Each request includes:
\begin{itemize}
    \item System prompt: concise reasoning format specification.
    \item User prompt: templated few-shot exemplars + target question.
    \item Decoding: \texttt{temperature=1.0}, \texttt{top\_p=1.0}.
    \item Timeout: 120\,s.
\end{itemize}

Tenacity applies exponential backoff (max 5 attempts) on transport errors.

\subsection{Formatting contract}
Required output:
\begin{verbatim}
Reasoning:
- step 1
- step 2
...
Final Answer: <number>
\end{verbatim}
No extra trailing text. Final line must start exactly with \texttt{Final Answer:}. Parser scans lines to extract the numeric suffix.

\subsection{Parsing and correctness}
\begin{itemize}
    \item Gold answer number: extracted from GSM8K \texttt{answer} field after the \texttt{####} delimiter using a final-number regex.
    \item Teacher answer number: last numeric token on the \texttt{Final Answer:} line (regex supports integers, decimals, scientific notation).
    \item Correctness: strict numeric equality (Python int/float comparison). Any parsing failure or missing final line marks the sample invalid (counted as failed).
\end{itemize}

\subsection{Recording schema}
Each raw sample appended as a JSONL record with fields:
\texttt{question\_id, sample\_id, question, gold\_answer\_text, gold\_answer\_number, teacher\_answer\_text, teacher\_answer\_number, is\_correct, finish\_reason, usage\{prompt\_tokens, completion\_tokens\}, latency\_ms, model, request\_id}. Latency is wall-clock from request dispatch to response receipt.

\subsection{Resumability}
Before generating new samples, existing JSONL is scanned to build a per-question set of occupied \texttt{sample\_id}s. Missing IDs in \([0, N)\) are generated where completed questions are skipped entirely.

\subsection{Progress and metrics}
A live tqdm progress bar displays:
\begin{itemize}
    \item \textbf{cost}: computed continuously as
    \[
      \text{cost} = \frac{T_{in}}{10^6} c_{in} + \frac{T_{out}}{10^6} c_{out}
    \]
    with $c_{in}=c_{out}=\$0.27$/M tokens (taken from OpenRouter).
    \item \textbf{tokens\_in / tokens\_out}: cumulative prompt / completion tokens.
    \item \textbf{acc}: running fraction of \texttt{correct\_samples / total\_samples}.
    \item \textbf{ok/err}: successful vs failed attempts.
\end{itemize}

\subsection{Failure handling}
Failures include HTTP errors, timeouts, parsing errors, or malformed responses. Each increments \texttt{failed\_samples} and no partial record is written. Retries only cover transport exceptions—logical format violations are not retried to avoid cost inflation.

\subsection{Quality outcome}
Oversampling plus strict filtering yields a pool of diverse, correctly formatted rationales per question (typically multiple distinct reasoning paths). This supports effective downstream distillation while keeping total spend within the budget envelope.

\section{Cleaning}
This stage removes incorrect, duplicate, and extreme-length samples to yield a high-signal supervision set.

\subsection{Input}
Raw JSONL: \texttt{artifacts/data/raw/dataset.jsonl} containing 30{,}000 generations (1{,}000 questions $\times$ 30 samples), each with correctness flag, usage tokens, and rationale text.

\subsection{Operations}
\begin{enumerate}
    \item \textbf{Load \& index}: Read JSONL into a DataFrame indexed by \texttt{(question\_id, sample\_id)} for deterministic ordering.
    \item \textbf{Correctness filter}: Keep only rows where \texttt{is\_correct = True}. Removes all teacher numeric mismatches; prevents label noise (primary gain).
    \item \textbf{Exact rationale dedup}: Drop duplicates on \texttt{teacher\_answer\_text}. Prevents overweighting identical reasoning paths.
    \item \textbf{Token-length outlier removal}: Compute per-row \texttt{total\_tokens = prompt\_tokens + completion\_tokens}. Remove rows with \texttt{total\_tokens} $\ge$ 99th percentile (threshold 1392 tokens). This trims extreme verbose rationales that inflate sequence length and memory without adding diversity.
    \item \textbf{Persist}: Reset index and write cleaned JSONL to \texttt{artifacts/data/cleaned/dataset.jsonl}.
\end{enumerate}

\subsection{Justification}
\begin{itemize}
    \item \textbf{Strict correctness}: Ensures the student never trains on wrong final answers (no noisy label memorization).
    \item \textbf{Deduplication}: Reduces redundant gradient steps; improves effective rationale variety per batch.
    \item \textbf{Length capping}: Controls worst-case sequence length to avoid GPU OOM and keeps average tokens closer to the median (739 total tokens).
\end{itemize}

\subsection{Outcome}
Resulting size: 24{,}952 cleaned samples (from 30k raw). Teacher correctness in retained set is 100\%. Distribution tails (>\,99th percentile) removed while preserving all final answer lines.

\section{Preprocessing}
Transforms the cleaned corpus into the minimal supervision schema consumed by training scripts.

\subsection{Input}
Cleaned JSONL: \texttt{artifacts/data/cleaned/dataset.jsonl} (indexed by \texttt{question\_id}, \texttt{sample\_id}).

\subsection{Output}
Select only: \texttt{question}, \texttt{gold\_answer\_text}, \texttt{gold\_answer\_number}, \texttt{teacher\_answer\_text} \quad (rationale + final answer line), \texttt{teacher\_answer\_number}
Index columns are preserved until final write, then flattened.

\subsection{Artifact}
Write to \texttt{artifacts/data/processed/dataset.jsonl}.

\subsection{Schema Rationale}
Minimal fields reduce I/O and memory footprint, remove unused metadata (latency, token usage) for training, and keep both textual and numeric answers to enable strict evaluation parsing and auxiliary checks.

\subsection{Integrity}
All records already passed strict correctness; no further validation required. Numeric fields (\texttt{gold\_answer\_number}, \texttt{teacher\_answer\_number}) support exact-match evaluation; textual rationale preserved verbatim for supervised reasoning.

\section{Prompt templates}\label{sec:prompt_templates}
This section touches on the fixed prompt templates used across generation, training, and benchmarking. Templates are versioned as plain text under \texttt{code/prompts/}. No runtime mutation other than \texttt{\{question\}} placeholder replacement.

\subsection{Files and roles}
\begin{itemize}
    \item \textbf{CoT system} (\texttt{code/prompts/cot/system.txt}): Instructs concise multi-step reasoning (2--8 steps) plus a terminating final answer line.
    \item \textbf{CoT user} (\texttt{code/prompts/cot/user.txt}): Provides three few-shot exemplars (question, reasoning steps, final answer) followed by the target question placeholder.
    \item \textbf{Label-only system} (\texttt{code/prompts/label\_only/system.txt}): Constrains output to a single final answer line with no reasoning.
    \item \textbf{Label-only user} (\texttt{code/prompts/label\_only/user.txt}): Supplies same three few-shot answer-only exemplars plus the target question placeholder.
\end{itemize}

\subsection{Formatting contracts}
\begin{itemize}
    \item \textbf{CoT prompts}: 
    \begin{verbatim}
Reasoning:
- step 1
- step 2
...
Final Answer: <number>
    \end{verbatim}
    2--8 dash-prefixed steps; arithmetic shown explicitly; no trailing text after the final answer line.
    \item \textbf{Label-only prompts}: Single line \texttt{Final Answer: <number>} No units, commentary, or extra lines.
\end{itemize}

\section{Training}
Goal of training is to fine-tune a single 3B decoder-only model (Qwen/Qwen2.5-3B) with QLoRA under 4-bit NF4 quantization and LoRA adapters, using bfloat16 precision.

\subsection{Model, quantization, and PEFT}
\begin{itemize}
    \item Base: \texttt{Qwen/Qwen2.5-3B} loaded with BitsAndBytes 4-bit (\texttt{nf4}, double-quantization), compute dtype bf16; \texttt{device\_map="auto"}.
    \item Gradient checkpointing enabled; \texttt{model.config.use\_cache = false} during training.
    \item Tokenizer: fast; if \texttt{pad\_token} is missing, set to \texttt{eos\_token} and propagate pad IDs to model and generation config.
    \item LoRA: \texttt{r=16}, \texttt{alpha=32}, \texttt{dropout=0.05}, \texttt{bias="none"} on \{\texttt{q\_proj}, \texttt{k\_proj}, \texttt{v\_proj}, \texttt{o\_proj}, \texttt{gate\_proj}, \texttt{up\_proj}, \texttt{down\_proj}\}; task type \texttt{CAUSAL\_LM}.
\end{itemize}

\subsection{Data split and reproducibility}
\begin{itemize}
    \item Input: \texttt{artifacts/data/processed/dataset.jsonl}.
    \item Split: 80/20 by \texttt{question\_id} (\texttt{TRAIN\_SPLIT=0.8}) with \texttt{seed=42} ensures no question leakage across splits.
    \item Label-only mode dedup: keep a single record per \texttt{question\_id} (first by \texttt{sample\_id}) to avoid multiple labels per question.
    \item Seeds for Python, PyTorch, CUDA, and accelerate set to 42; TF32 enabled on matmul/CuDNN.
\end{itemize}

\subsection{Prompting and target formatting}
\begin{itemize}
    \item Prompts: training/inference use the same templates.
    \begin{itemize}
        \item SCoTD: \texttt{prompts/cot/system.txt} + \texttt{prompts/cot/user.txt}.
        \item Label-only: \texttt{prompts/label\_only/system.txt} + \texttt{prompts/label\_only/user.txt}.
    \end{itemize}
    \item Targets:
    \begin{itemize}
        \item SCoTD: full teacher rationale text including the final line.
        \item Label-only: a single line "Final Answer: \textless number\textgreater" built from the gold number.
    \end{itemize}
\end{itemize}

\subsection{Sequence construction and truncation}
\begin{itemize}
    \item Max sequence length: 2048.
    \item Encoding: concatenate prompt IDs with answer IDs; only answer tokens contribute to loss.
    \item Truncation policy: always preserve the full answer. If needed, trim the prompt from the left so that $\lvert\text{prompt}\rvert + \lvert\text{answer}\rvert \le 2048$.
    \item Loss mask: labels are \texttt{-100} for all prompt positions and the token IDs of the answer (incl. EOS) for answer positions.
\end{itemize}

\subsection{Batching and collator}
\begin{itemize}
    \item Dynamic padding to the longest in batch; pad to a multiple of 8 for Tensor Cores; attention mask computed as \texttt{input\_ids != pad\_id}.
    \item \texttt{group\_by\_length=true} to reduce padding waste.
\end{itemize}

\subsection{Optimization and schedule}
\begin{itemize}
    \item Trainer: HuggingFace \texttt{Trainer}.
    \item Optimizer: \texttt{paged\_adamw\_8bit}.
    \item Scheduler: linear; \texttt{warmup\_ratio} per mode (below).
    \item Precision: bf16 if hardware supports it, fp16 otherwise.
    \item Dataloading: \texttt{num\_workers=4}, \texttt{pin\_memory=true}.
    \item Checkpointing: \texttt{save\_strategy="steps"}; \texttt{load\_best\_model\_at\_end=true} with \texttt{metric\_for\_best\_model="eval\_loss"} and \texttt{greater\_is\_better=false}.
    \item Logging: TensorBoard, \texttt{logging\_steps=10}.
\end{itemize}

\subsection{Mode-specific hyperparameters}
\noindent SCoTD (rationale+answer):
\begin{itemize}
    \item epochs=10, learning rate=$2{\times}10^{-5}$.
    \item per-device train/eval batch size=8, grad accumulation=2.
    \item warmup\_ratio=0.03, weight\_decay=0.
    \item save/eval every 200 steps.
\end{itemize}
\noindent Label-only (answer line only):
\begin{itemize}
    \item epochs=20, learning rate=$1{\times}10^{-5}$.
    \item per-device train/eval batch size=16, grad accumulation=1.
    \item warmup\_ratio=0.1, weight\_decay=0.1, max\_grad\_norm=1.0.
    \item save/eval every 50 steps.
\end{itemize}

\subsection{Export and sanity check}
\begin{itemize}
    \item After training, reload the frozen 4-bit base and attach the best LoRA checkpoint (\texttt{PeftModel.from\_pretrained}) for inference.
    \item Deterministic, greedy generation (\texttt{do\_sample=false}) used for quick sanity checks on held-out questions; decoding respects set \texttt{pad\_token\_id} and \texttt{eos\_token\_id}.
\end{itemize}

\section{Benchmarking}
Benchmarking evaluates the base model and all student checkpoints on the GSM8K test split under strict exact-match parsing of the final numeric answer.

\subsection{Objective}
Measure accuracy ($\#correct \div \#total$) for:
\begin{itemize}
    \item Base model with CoT prompting.
    \item Base model with label-only prompting.
    \item All SCoTD LoRA checkpoints.
    \item All label-only LoRA checkpoints.
\end{itemize}
All generations use deterministic greedy decoding (no sampling, no self-consistency).

\subsection{Dataset and gold answers}
The GSM8K test split is loaded via \texttt{load\_dataset(GSM8K\_PATH, name="main", split="test")}. Gold numeric answers are extracted from the original \texttt{answer} field using \texttt{parse\_gold\_answer\_number}, which regex-parses the final integer after ''\#\#\#\#'' delimiter.

\subsection{Prompt construction}
Two builders create inference prompts:
\begin{itemize}
    \item \texttt{build\_prompt\_cot}: concatenates the CoT system template and user template with the target question; expects rationale steps then a final line.
    \item \texttt{build\_prompt\_label\_only}: analogous, but the system template instructs answer-only output.
\end{itemize}
Both yield a single plain-text prompt string; any truncation (left side) preserves the answer region.

\subsection{Model loading and quantization}
Each run loads either the base model (\texttt{Qwen/Qwen2.5-3B}) or a LoRA checkpoint path. All models use 4-bit NF4 quantization (\texttt{BitsAndBytesConfig} with double quantization) and bf16 compute. \texttt{device\_map="auto"}. \texttt{model.config.use\_cache = True} activates KV caching for faster greedy decoding.

\subsection{Tokenizer settings}
Tokenizer is fast, left-padding and left-truncating. Pad token ID is propagated to \texttt{generation\_config}. Batches are tokenized with dynamic padding; no special postprocessing beyond stripping decoded prompt prefixes from responses.

\subsection{Generation configuration}
\begin{itemize}
    \item Greedy: \texttt{do\_sample=False}.
    \item \texttt{max\_new\_tokens=1024} to allow long rationales.
    \item Cache usage: \texttt{use\_cache=True}.
    \item Precision aids: TF32 enabled on matmul/CuDNN for potential minor throughput gains.
\end{itemize}

\subsection{Batching}
Batch size reflects output length expectations:
\begin{itemize}
    \item Label-only: 128 (short single-line answers).
    \item CoT: 16 (long multi-step rationales).
\end{itemize}
Processing iterates over slices of question lists; each batch calls \texttt{model.generate} once.

\subsection{Parsing and strict evaluation}
Predicted texts are trimmed by removing the prompt prefix and then parsed with \texttt{parse\_teacher\_final\_answer}, which searches for a final line starting with ``Final Answer:'' and extracts the trailing numeric token. Failures (missing line, malformed number, regex miss) raise \texttt{ParsingError}, such cases count as incorrect. This strictness penalizes:
\begin{itemize}
    \item Format deviations (e.g., extra commentary after the answer line).
    \item Missing prefix ``Final Answer:''.
    \item Multiple numeric candidates (the parser only accepts the final numeric token on that line).
\end{itemize}
Accuracy is the fraction of predictions whose parsed number equals the gold number (exact equality). No normalization (e.g., stripping commas or units) is applied.

\subsection{Checkpoint iteration and resource cleanup}
A \texttt{RUNS} list enumerates all label-only checkpoints (sorted by numeric suffix), all SCoTD checkpoints, then the two base-mode baselines. After each run:
\begin{itemize}
    \item Metrics (accuracy, count) are appended to a results list.
    \item Predictions DataFrame rows receive the model name.
    \item \texttt{cleanup\_model} deletes model/tokenizer objects, triggers garbage collection, empties CUDA cache, and synchronizes to reduce fragmentation before loading the next checkpoint.
\end{itemize}
This prevents cumulative GPU memory growth across many sequential loads.

\subsection{Artifacts}
All per-question predictions across runs are concatenated and written to \texttt{predictions.csv} with columns: \texttt{model} (run identifier), \texttt{question}, \texttt{prediction} (raw decoded answer text), \texttt{gold\_answer} (numeric gold).

\section{Environment setup}
Fine-tuning and benchmarking runs were executed on RunPod node with the following specs:

\subsection{Hardware}
\begin{itemize}
    \item Single GPU: RunPod RTX 4090 (24GB VRAM), RAM 64GB.
    \item Pod cost: \$0.34/hour.
\end{itemize}

\subsection{Software}
\begin{itemize}
    \item Python 3.11; package manager: \texttt{uv}.
    \item CUDA 12.9; PyTorch 2.8.0.
    \item Env: \texttt{TOKENIZERS\_PARALLELISM=false}.
\end{itemize}